\section{Glossary}\label{app:glossary}
\begin{description}[leftmargin=2.2em,style=nextline,itemsep=1pt,font=\normalfont\bfseries]
\small
\item[Query, key, value ($Q$, $K$, $V$)] The three vectors each token contributes to attention. A query asks, keys are matched against it, and values are what gets averaged.
\item[Head; MHA; GQA] An independent copy of the attention operation. Multi-head attention gives every head its own keys and values; grouped-query attention lets $G$ query heads share one key/value head.
\item[Softmax] The map from scores to positive weights that sum to one, $p_j=e^{s_j}/\sum_{j'}e^{s_{j'}}$.
\item[Online softmax] Maintaining $(m,\ell,a)$ so that keys can be processed block by block in any order (\cref{eq:online}).
\item[Causal mask; window; sink] Causal: a query may only use keys at earlier or equal positions. Window: only the last $w$ keys. Sink: the first few tokens of the sequence, which many heads attend to regardless of content.
\item[RoPE] Rotary position embedding: the model rotates queries and keys by an angle that depends on token position before the dot product. The extracted real tensors already include it.
\item[Rank; shard; rows per rank $c$] An MPI process; the part of the sequence it owns; $c=N/P$.
\item[Contiguous / zigzag partition] One run of positions per rank, or one early plus one late chunk per rank.
\item[All-gather; ring] Copy every shard to every rank at once; or pass shards from neighbour to neighbour so that each visits every rank.
\item[Two-sided / one-sided communication (RMA)] Send/receive pairs in which both ranks call MPI; or direct access to another rank's exposed memory.
\item[Window; epoch; passive / active target; fence] The exposed memory region; the period in which access is legal; whether the target takes part in synchronisation; the collective call that ends an active-target epoch.
\item[Progress] The internal work an MPI library must do to move data; some libraries only do it inside MPI calls.
\item[Control M0; $E_0$; $E_\text{comm}$] A run that gathers all keys and values before the clock starts and then computes without calling MPI; its efficiency $T_S/(P\,T_\text{M0})$; the speed of a variant relative to it, $T_\text{M0}/T$.
\item[Host-staged / CUDA-aware MPI] Exchange host copies of the blocks with MPI and upload each received block to the GPU, or hand device pointers to the MPI library.
\item[Tile; packing] A small block of query rows or keys that fits in cache; copying a key tile into a layout that suits the inner loop.
\item[Useful FLOPs; GFLOP/s; TFLOP/s] $2(d+d_v)$ floating-point operations per allowed query--key pair and head; $10^9$ or $10^{12}$ of them per second.
\item[Clock; turbo] The frequency at which a core executes instructions. The processor raises it above the base value under load (turbo) and lowers it again when many cores are active.
\item[Ping-pong] A microbenchmark in which one rank sends a message to another and receives it back; half the round trip is the time per message.
\item[Warm-up; repetition] An untimed execution of a configuration before it is measured; one timed execution of it.
\item[Speedup; efficiency] $S=T_S/T_P$ with $T_S$ the best single-core time; $E=S/\text{cores}$.
\item[Strong / weak scaling] More cores for the same problem; or a problem that grows with the cores.
\end{description}

\section{Environment and reproduction}\label{app:environment}
\begin{table}[!ht]
\centering\small
\begin{tabularx}{\textwidth}{@{}l Y@{}}
\toprule
Item & Value (recorded by every job in its \code{environment.txt})\\
\midrule
Cluster & NUS Vanda; all results in this report\\
CPU nodes & 2 $\times$ Intel Xeon Platinum 8452Y (Sapphire Rapids), 36 cores per socket, 72 cores per node, no hyper-threading; 48\,KiB L1d and 2\,MiB L2 per core, 67.5\,MiB L3 per socket; 512\,GiB memory; two NUMA domains\\
Clock & 2.0\,GHz base; \n{hw.freq.1}\,GHz sampled with up to nine active cores of a socket, \n{hw.freq.72}\,GHz with all 72 active (\cref{fig:hardware})\\
Network & shared memory inside a node; between nodes TCP over Ethernet (Open MPI selected the \code{ob1} messaging layer with the \code{sm} and \code{tcp} transports; the UCX layer did not initialise on the nodes' RDMA-capable Ethernet adapters)\\
CPU software & GCC 14.3.0, \code{-O3 -march=native} (\code{sapphirerapids}), \code{-mprefer-vector-width=512} for the compute core, all tensors and scratch buffers aligned to 4096 bytes; Open MPI 5.0.8 (PRRTE 3.0.11, UCX 1.19.0); OpenMP 4.5 (libgomp); CMake 3.31.8\\
GPU nodes & 2 $\times$ NVIDIA A40 (46\,GiB, compute capability 8.6) per node, driver 575.57.08; a GPU job receives 36 cores per GPU, of which the benchmark uses one thread per rank\\
GPU software & NVIDIA HPC SDK 26.1: \code{nvcc} 12.9, HPC-X 2.25.1 (CUDA-aware Open MPI 4.1.9, UCX); host code compiled with GCC 11.5. This MPI has no PBS support: on two nodes the second rank is started over ssh from the node file\\
Tile sizes & $B_q=\n{tiles.used.bq}$ query rows, $B_k=\n{tiles.used.bk}$ keys (\cref{fig:tiles}); a rank with $T>1$ threads uses query tiles of at most $\lceil c/(4T)\rceil$ rows (\cref{sec:xdesign})\\
Binding & one core per rank (\code{--map-by slot:PE=$T$ --bind-to core}); OpenMP threads pinned to cores (\code{OMP\_PLACES=cores}, \code{OMP\_PROC\_BIND=close}); ranks fill socket 0, then socket 1, then the next node\\
Protocol, CPU timings & untimed warm-up of at least one repetition and at least five seconds, then three measured repetitions in a row; time of a repetition = maximum over ranks; reported value = median of the three; each repetition validated on sampled rows\\
Protocol, GPU timings & one untimed warm-up repetition (no minimum time), then three measured repetitions, each validated on sampled rows; the time of a repetition includes the upload of the first block and the download of the result\\
Other protocols & peak memory: one run per configuration, no warm-up, no validation copy. Clock ramp: 400 consecutive repetitions without warm-up. Hardware counters: fixed work under \code{perf stat} (\cref{tab:align}, without validation; \cref{fig:hardware}b,c, a fixed number of validated repetitions per rank count). Network and ring-exchange microbenchmarks: their own repetition schemes. Unit tests: every row compared\\
Real data & Qwen2.5-0.5B (Hugging Face), PyTorch 2.5.0 and transformers 4.47.1 for the extraction only; text: Project Gutenberg \#2600. The SHA-256 digests of the model files, the text and the extracted tensors are kept with the results\\
\bottomrule
\end{tabularx}
\caption{Hardware, software and measurement protocol.}
\label{tab:environment}
\end{table}

\begin{sloppypar}
\paragraph{Source code.} The code of this project, with a README that explains how
to build and run it, is in the GitHub repository
\url{https://github.com/Zepeng-Chen/ring-attention-mpi}, which is public from the
submission deadline on. No BLAS, LAPACK or other
numerical library is used; the attention kernels, the MPI strategies and the CUDA
kernels are written out in \code{src/} and \code{include/attn/}.

\paragraph{Reproduction.} The README lists the build and run commands. Each experiment
of this report is one suite of \code{scripts/hpc/jobs/experiment.pbs} (CPU),
\code{gpu\_experiment.pbs} (GPU), \code{real\_data.pbs} or \code{unit\_tests.pbs};
\code{scripts/hpc/cluster.sh} stages an immutable copy of the source, named by a hash
of exactly the staged files and checked against that hash after copying, submits
the job and retrieves the result directory. \code{scripts/build\_report.sh} deletes the generated figures,
tables and numbers, rebuilds them and the PDF from the retrieved files, and stops if
a number has no data, a figure is missing or a reference is undefined.

\paragraph{One audited set of records.} Figures, tables and the measured values in
the text are computed by three scripts from the same \n{audit.runs} run directories
(\n{audit.files} result files). The directories are named explicitly in
\code{results/selected\_runs.json}, together with the \n{audit.excluded} run
directories that are \emph{not} used and the reason for each exclusion (measured
with an earlier protocol or source, diagnostics, jobs that failed or never started);
no script chooses a run by itself. Timings, throughputs, errors and statistics are
never typed into the text; what is typed are descriptions of the hardware, of the
inputs and of the experimental design. \code{scripts/audit\_results.py} checks the
selected records before every build and stops it if a rule is broken: every run
ended normally on Vanda; a validated repetition may fail only in the declared
stress study (\n{audit.failed} of \n{audit.validated}); every timed configuration
has its three repetitions with finite times, and CPU runs their five-second warm-up
(the declared exceptions are those of \cref{tab:environment}); all CPU jobs were
built from byte-identical compute and benchmark-driver sources and all GPU jobs
from identical kernel sources, and the submitted compute and kernel sources are
those files; the build flags and tile sizes are the stated ones; the unit tests ran
on Vanda on exactly the submitted files. Two programs were changed after their
records were computed, both as a result of a code review. The benchmark driver now
rejects infinite values of two options (the warm-up time and the tolerance scale).
The audit verifies that undoing those two conditions restores the measured file
byte for byte, and that none of the \n{audit.driver.commands} recorded command
lines gave either option such a value, so every recorded run is read the same way
by both versions. The analysis program of \cref{sec:real} now handles two edge
cases (a window whose weights all underflow, an exact output of zero). The audit
checks on the stored records that neither case occurs in them
(\n{audit.e11.records} records, smallest window probability
$\n{audit.e11.minmass}$, no zero output), in which case the old and the new program
compute the same values. The audit also records which result
files each figure, table and group of numbers was computed from.
\Cref{tab:provenance} is that record, by job.
\end{sloppypar}

\inputifexists{report/generated/table_provenance.tex}

